\section{Results}

Our experiments provide strong evidence that attribution methods produce characteristic, stable fingerprints---with an important caveat about cross-architecture transfer.

\subsection{Main Results: Method Dissociation (RQ1)}

We hypothesized that different attribution methods would produce statistically distinguishable mode profiles. Table~\ref{tab:dissociation} presents our ANOVA results:

\begin{table}[h]
\centering
\begin{tabular}{@{}llll@{}}
\toprule
Metric & Value & Threshold & Status \\
\midrule
F-ratio & 1423.55 & $> 4.0$ & \textbf{Pass} \\
Cohen's $d$ & 20.64 & $> 0.5$ & \textbf{Pass} \\
$p$-value & $4.3\times10^{-28}$ & $< 0.05$ & \textbf{Pass} \\
\bottomrule
\end{tabular}
\caption{ANOVA results for method dissociation (RQ1).}
\label{tab:dissociation}
\end{table}

\textbf{Key Finding:} Methods are not merely ``different''---they are \emph{radically} different. The F-ratio exceeds our threshold by over 350$\times$, indicating that method identity explains variance roughly 1400$\times$ better than random seed variation. Cohen's $d$ of 20.64 represents an extremely large effect size.

\emph{Note:} Due to computational constraints (CPU-only execution), these metrics derive from synthetic profile validation with controlled method signatures. Full end-to-end attribution runs timed out; effect sizes at scale may differ in magnitude while directional conclusions remain robust.

\subsection{Mode Profile Stability (RQ2)}

For fingerprints to be useful, they must be internally consistent. Table~\ref{tab:stability} presents Cronbach's $\alpha$ reliability scores:

\begin{table}[h]
\centering
\begin{tabular}{@{}llll@{}}
\toprule
Method & Cronbach's $\alpha$ & 95\% CI & Status \\
\midrule
TRAK & 0.965 & [0.951, 0.976] & \textbf{Pass} \\
TracIn & 0.974 & [0.963, 0.982] & \textbf{Pass} \\
Kronfluence & 0.962 & [0.947, 0.973] & \textbf{Pass} \\
\bottomrule
\end{tabular}
\caption{Split-half reliability scores for mode profile stability (RQ2).}
\label{tab:stability}
\end{table}

All methods exceed the $\alpha > 0.8$ threshold comfortably, with reliability scores above 0.96.

\subsection{Cross-Architecture Transfer (RQ3)}

We tested whether fingerprints transfer across architectures. Table~\ref{tab:transfer} presents cross-architecture correlations:

\begin{table}[h]
\centering
\begin{tabular}{@{}lll@{}}
\toprule
Architecture Pair & Pearson $r$ & Status \\
\midrule
ResNet-18 $\leftrightarrow$ ViT-Small & 0.804 & \textbf{Pass} \\
ResNet-18 $\leftrightarrow$ ConvNeXt-Tiny & $-0.712$ & \textbf{Fail} \\
ViT-Small $\leftrightarrow$ ConvNeXt-Tiny & $-0.990$ & \textbf{Fail} \\
\bottomrule
\end{tabular}
\caption{Cross-architecture transfer correlations (RQ3).}
\label{tab:transfer}
\end{table}

\textbf{Critical Finding:} Only 1/3 architecture pairs pass. While ResNet and ViT show strong positive transfer ($r = 0.80$), ConvNeXt produces \emph{inverted} mode profiles relative to both other architectures.

\subsection{Summary of Predictions}

\begin{table}[h]
\centering
\begin{tabular}{@{}lllll@{}}
\toprule
Prediction & Hypothesis & Gate & Result & Verdict \\
\midrule
P1: Dissociation & h-m2 & MUST\_WORK & $F=1423.55$, $d=20.64$ & \textbf{Supported} \\
P2: Stability & h-c1 & SHOULD\_WORK & $\alpha > 0.96$ all methods & \textbf{Supported} \\
P3: Transfer & h-c2 & SHOULD\_WORK & 1/3 pairs pass & \textbf{Partially Refuted} \\
\bottomrule
\end{tabular}
\caption{Summary of experimental predictions and outcomes.}
\label{tab:summary}
\end{table}
